%=====================================================================
\chapter{Hybrid Language Model and Indicator Framework}
\label{ch:llm}
%=====================================================================

\textbf{Research Objective 2.} \emph{To design a hybrid framework that combines large language
models, sentiment analysis techniques and conventional financial indicators, focusing on
optimising deep learning architectures for real-time stock market predictions.}

\vspace{0.4em}
Chapters~\ref{ch:fusion} and~\ref{ch:dynfusion} maximised the number of sources and showed that
combining them helps. This chapter does the opposite. It isolates a single source --- news text
read by a transformer language model --- and asks precisely what that source adds to a
conventional indicator set at the horizon and granularity at which a practitioner actually
takes a position: the next-day direction of an individual equity.

The answer is more interesting than a headline accuracy figure, and it is reported in full
including the part that does not flatter the method. Language-model sentiment turns out to be
the single most important feature in the fused space and to carry information orthogonal to
price, while delivering only a small increment in raw directional accuracy, with every model
examined operating close to the random classifier. Section~\ref{sec:llm-limit} treats that
combination as the finding rather than as a disappointment, and Section~\ref{sec:llm-implication}
explains why it redirects the remainder of the thesis toward the decision layer.

The work reported in this chapter appears in publication~\pubcite{P3}; the architectural
feasibility evidence in Section~\ref{sec:llm-feasibility} draws on~\pubcite{P1}.

\section{Design Rationale}
\label{sec:llm-rationale}

\subsection{Why isolate a single source}

A multi-source architecture that outperforms its baselines has demonstrated that the collection
of sources is useful. It has not established what any individual source contributes, because the
sources are confounded: any gain could come from any of them, or from their interaction, or from
the additional model capacity the architecture required to accommodate them.
Section~\ref{sec:msfn-ablation} addressed part of this problem by ablation, but ablation within a
fused model measures the \emph{marginal} contribution of a stream in the presence of the others,
which is not the same as its standalone informativeness.

This chapter therefore constructs the cleanest available comparison: one textual source, one
conventional feature set, one fusion step, and baselines that differ from the proposed model in
exactly that one source. Every model sees the same sessions, the same partition, the same
protocol. Whatever difference appears is attributable to the sentiment feature.

\subsection{Why a continuous score rather than a label}

The sentiment encoder emits a continuous polarity score $s_t \in [-1,+1]$ together with a
confidence, rather than a categorical positive, neutral or negative label. The reason is that
intensity is informative and categorisation destroys it. A routine quarterly result described in
mildly favourable terms and a transformative regulatory approval both map to ``positive''; the
market does not treat them alike. Retaining the continuous value preserves the distinction for
the downstream model, and it also makes the empirical distribution of sentiment available as a
diagnostic, which Section~\ref{sec:llm-distribution} uses.

\subsection{Why a tree ensemble rather than a deep network}

The model operating on the fused feature space is a gradient-boosted tree
ensemble~\cite{chen2016xgboost} rather than a deep network, for three reasons stated in
Section~2.6.3 and restated here because the choice is deliberate. Tree ensembles are invariant to
monotone rescaling, which matters when combining a bounded sentiment score with unbounded volume
accumulators. They handle heterogeneous feature types without embedding layers. And their
split-gain statistics supply a feature-importance measure intrinsic to the fitted model rather
than a post-hoc approximation --- which is what makes the central finding of
Section~\ref{sec:llm-importance} a measurement rather than an estimate.

There is also a practical reason. A design goal of this study was that the entire pipeline should
run to completion in a freely accessible hosted notebook environment, so that the result can be
independently re-derived by a reader without specialised hardware. A tree ensemble over a
few thousand daily observations satisfies that constraint; a large recurrent architecture over
the same data does not, and would in any case be over-parameterised for the sample size.

\section{Framework Architecture}
\label{sec:llm-arch}

The framework is organised as five interconnected layers, shown in Figure~\ref{fig:llm}: data
collection, technical indicator computation, sentiment analysis, feature fusion, and model
training with evaluation.

\begin{figure}[!htb]
\centering
\scalebox{0.90}{\begin{tikzpicture}[node distance=3mm]
\node[blkd,minimum width=32mm,minimum height=8mm] (a1) {Historical prices\\daily OHLCV};
\node[blkd,right=10mm of a1,minimum width=32mm,minimum height=8mm] (a2) {Financial news\\headlines and articles};

\node[blk,below=7mm of a1,minimum width=32mm,minimum height=11mm] (b1)
  {\textbf{Technical layer}\\15 indicators: SMA, RSI,\\MACD, Bollinger, VWAP};
\node[blk,below=7mm of a2,minimum width=32mm,minimum height=11mm] (b2)
  {\textbf{Sentiment layer}\\transformer encoder $\rightarrow$\\continuous score $s_t \in [-1,+1]$};

\node[blkg,below=7mm of b1,minimum width=32mm,minimum height=8mm] (c1)
  {$z$-score standardisation};
\node[blkg,below=7mm of b2,minimum width=32mm,minimum height=8mm] (c2)
  {min--max scaling and\\session alignment};

\node[blka,below=7mm of c1,minimum width=74mm,minimum height=8mm,xshift=21mm] (d)
  {\textbf{Feature fusion layer} --- early, feature-level concatenation};
\node[blka,below=5mm of d,minimum width=74mm,minimum height=8mm] (e)
  {\textbf{Model layer} --- gradient-boosted ensemble on the fused space};
\node[blkd,below=5mm of e,minimum width=74mm,minimum height=10mm] (f)
  {\textbf{Evaluation layer} --- temporal cross-validation, directional accuracy,\\
   Diebold--Mariano testing, confidence calibration, simulated trading};

\foreach \a/\b in {a1/b1,a2/b2,b1/c1,b2/c2} \draw[ar] (\a) -- (\b);
\draw[ar] (c1.south) -- ++(0,-2mm) -| (d.north);
\draw[ar] (c2.south) -- ++(0,-2mm) -| (d.north);
\draw[ar] (d) -- (e);
\draw[ar] (e) -- (f);
\node[lbl,right=3mm of b2,align=left,xshift=2mm] {contextual reading\\replaces lexicon\\polarity};
\end{tikzpicture}}
\caption[Hybrid language model and technical indicator framework]{Hybrid framework combining
large-language-model sentiment extraction with conventional financial indicators. The continuous
sentiment score preserves intensity that categorical classification would discard.}
\label{fig:llm}
\end{figure}

\subsection{Data collection layer}

Market data are retrieved for Indian large-capitalisation equities --- TCS, RELIANCE and
HDFCBANK --- as daily open, high, low, close and volume records over a multi-year period
sufficient to span several distinct market conditions. The three names were chosen to cover
different sectors (information technology, energy and conglomerate, banking) and because all
three carry news coverage dense enough to produce a meaningful daily sentiment aggregate; a
thinly covered name would yield a sentiment series that is mostly missing, which tests the
imputation rule rather than the hypothesis.

News acquisition uses a two-stage strategy combining syndicated news feeds with targeted
retrieval, and applies filtering at two levels: non-financial content is excluded, and sources
are prioritised by established credibility in financial reporting. Items are validated for
temporal consistency --- an item whose time stamp is missing or implausible is discarded rather
than imputed --- and for content quality.

\subsection{Technical indicator layer}

Fifteen indicators are computed, following the definitions of Section~2.2. Trend is represented
by simple moving averages over 5, 10 and 20 sessions; momentum by the relative strength
index~\cite{wilder1978} and the moving average convergence divergence with its signal line;
volatility by Bollinger band measures and historical volatility; and participation by
volume-weighted average price derivatives and volume momentum oscillators. Each indicator was
selected for established efficacy in the financial literature and for computational efficiency,
the two criteria being in tension only rarely at this scale.

\subsection{Sentiment analysis engine}

Each headline is scored by a transformer sentiment encoder operating on the raw text. The
pipeline has four stages: text normalisation to remove markup and boilerplate; batched inference
for computational efficiency; extraction of a continuous polarity score and an associated
confidence; and daily aggregation over all items attributed to the session.

The contextual reading performed by a transformer is what distinguishes this stage from a lexicon
baseline. A lexicon scores ``the bank wrote back provisions after the regulator eased norms'' by
summing token polarities, several of which are negative in general usage; a contextual encoder
reads the sentence. This is the improvement in \emph{measurement} that the literature surveyed in
Section~3.2 establishes, and the entire purpose of this chapter is to determine how much of it
survives into forecast quality.

\subsection{Feature fusion layer}

Fusion is performed at feature level, before modelling. Technical indicators are standardised by
$z$-score, which equalises scales across features measured in incompatible units --- a price
level, a bounded oscillator and a volume ratio have no common scale. The sentiment score is
min--max scaled, which preserves its interpretable range rather than mapping it to an arbitrary
standardised scale. The two blocks are concatenated into a single input vector.

Temporal alignment is applied here and is the most error-prone step in the framework. The rule,
as in Section~2.5, is that an item is attributed to the first trading session beginning strictly
after its publication time stamp. An article published at 18:30 on a Tuesday is attributed to
Wednesday's session, not to Tuesday's, which has already closed. Attributing it to Tuesday would
be a form of look-ahead that leaves no trace in the code and inflates every downstream metric.

\subsection{Model training and evaluation layers}

A gradient-boosted tree ensemble is trained on the fused space. Three baselines are trained on
technical indicators alone under an identical protocol: random forest~\cite{breiman2001}, support
vector machine~\cite{cortes1995svm} and logistic regression. The target is the sign of the
next-session return.

The evaluation layer applies a chronological 80/20 partition with temporal cross-validation,
reports accuracy, precision, recall and $F_1$, tests differences by the Diebold--Mariano
statistic~\cite{diebold1995}, examines confidence calibration, and converts predictive accuracy
into economic terms through a simulated trading strategy. The implementation stack is recorded
in Table~\ref{tab:llm-stack} for reproducibility.

\begin{table}[!htb]
\centering
\caption[Implementation stack for the hybrid framework]{Implementation stack for the hybrid
language-model and indicator framework}
\label{tab:llm-stack}
\tabfont
\begin{tabular}{@{}L{4.0cm}L{9.6cm}@{}}
\toprule
\textbf{Component} & \textbf{Implementation} \\
\midrule
Environment & Hosted notebook environment, freely accessible, no specialised hardware required \\
Data processing & \texttt{pandas}, \texttt{numpy} \\
Market data & \texttt{yfinance}, \texttt{pandas\_datareader} \\
Machine learning & \texttt{scikit-learn}, \texttt{xgboost} \\
Language model & \texttt{transformers}, \texttt{torch} \\
Visualisation & \texttt{matplotlib}, \texttt{seaborn}, \texttt{plotly} \\
Design principle & Modular separation of data processing, feature engineering, model training and
evaluation, so that each subsystem can be tested independently \\
\bottomrule
\end{tabular}
\end{table}

\section{Results: Classification Performance}
\label{sec:llm-results}

Table~\ref{tab:llmres} and Figure~\ref{fig:llmres} report the comparison. The result must be read
carefully, because the headline number and the useful number are not the same.

\begin{table}[!htb]
\centering
\caption[Directional classification with LLM-extracted sentiment]{Single-name directional
classification with and without language-model-extracted sentiment features}
\label{tab:llmres}
\tabfont
\begin{tabular}{@{}L{4.5cm}L{3.3cm}rrrr@{}}
\toprule
\textbf{Model} & \textbf{Feature set} & \textbf{Accuracy} & \textbf{Precision} &
\textbf{Recall} & \textbf{$F_1$} \\
\midrule
Random Forest & technical only & 0.4982 & 0.4345 & 0.6033 & 0.5052 \\
Support Vector Machine & technical only & 0.4246 & 0.4246 & 1.0000 & 0.5961 \\
Logistic Regression & technical only & 0.4702 & 0.4257 & 0.7107 & 0.5325 \\
XGBoost with sentiment & technical + LLM sentiment & 0.4772 & 0.4103 & 0.5289 & 0.4621 \\
\midrule
\multicolumn{2}{@{}l}{Mean of technical-only baselines} & 0.4643 & --- & --- & --- \\
\multicolumn{2}{@{}l}{Relative improvement of the proposed model} & \textbf{+2.77\%} & --- &
--- & --- \\
\bottomrule
\end{tabular}
\end{table}

\begin{figure}[!htb]
\centering
\includegraphics[width=0.90\linewidth]{figs/f_models.pdf}
\caption[Comparative classification performance with sentiment features]{Comparative
classification performance. All models operate close to the random classifier at a single-name
daily horizon, which is the correct baseline against which the sentiment increment should be
judged.}
\label{fig:llmres}
\end{figure}

The sentiment-augmented model attains 47.72 per cent accuracy against a baseline mean of 46.43
per cent, a relative improvement of 2.77 per cent. It does \emph{not} attain the highest raw
accuracy: random forest reaches 49.82 per cent. Reporting only the first of these numbers would
misrepresent the finding, so both are stated prominently.

The precision--recall columns are instructive and are the reason accuracy alone is an inadequate
summary. The support vector machine attains perfect recall of 1.0000 with precision of 0.4246,
which is exactly the signature of a degenerate classifier: it predicts the positive class
unconditionally, and its precision is therefore simply the base rate of upward sessions in the
sample. Its $F_1$ of 0.5961 is the highest in the table and means nothing at all. Any evaluation
that ranked models by $F_1$ alone would select the one model that has learned nothing.

What makes the result meaningful is the structure beneath the aggregate, and the next four
sections examine it.

\section{Feature Importance: Sentiment Dominates the Fused Space}
\label{sec:llm-importance}

In the fused model, the sentiment score is the highest-gain predictor, ranking above established
technical features including the twenty-session moving average and the moving average convergence
divergence. The measure is the intrinsic split-gain statistic of the tree ensemble --- the total
reduction in loss attributable to splits on that feature --- rather than a post-hoc attribution,
so it reports what the fitted model actually relied upon.

The interpretation deserves care. Feature importance measures how heavily a model \emph{uses} a
feature, not how much predictive accuracy that feature delivers; a model can rely on a feature and
still predict poorly if the feature is the best available in a fundamentally hard problem. That is
precisely the situation here, and recognising it is what makes the combination of findings in this
chapter coherent rather than contradictory: sentiment is the most useful thing available, and the
most useful thing available is not very useful at this horizon.

The finding nonetheless establishes something the literature surveyed in Section~3.2 asserts
without demonstrating: language-model sentiment is not a marginal addition that a model tolerates
alongside its real features. When it is present, it is the feature the model leans on most.

\section{Orthogonality: Sentiment Carries Information Price Does Not}
\label{sec:llm-orthogonality}

The correlation structure of the fused feature space settles the question of whether sentiment is
redundant. Two correlations are decisive.

The correlation between the sentiment score and the price level is 0.04 --- effectively zero. The
correlation between the sentiment score and momentum indicators, principally the moving average
convergence divergence, is 0.74 --- substantial.

Both numbers are necessary to the interpretation and neither alone is sufficient. Sentiment aligns
with \emph{momentum}, which is expected: news and price both respond to the same underlying
events, so a stream of favourable reports and a rising momentum reading are two views of one
process. If this correlation had been near zero, one would have to suspect that the sentiment
pipeline was measuring noise. Simultaneously, sentiment carries information essentially
independent of the \emph{price level}, which is the condition under which fusion can help at all.
A feature perfectly correlated with an existing feature adds nothing regardless of how predictive
it is; a feature uncorrelated with the existing set can add information even if its standalone
predictive power is modest.

This orthogonality is the structural justification for the entire multi-source programme of this
thesis, measured directly at the single-name level rather than inferred from an aggregate
improvement. It is also consistent with the superadditivity observed in
Section~\ref{sec:msfn-ablation}: sources that carry non-overlapping information should
substitute imperfectly for one another, and both measurements are detecting the same property from
different directions.

\section{Asymmetry Across Sentiment Regimes}
\label{sec:llm-regimes}

Accuracy measured separately under positive, neutral and negative sentiment shows two patterns.

First, accuracy is higher under positive sentiment for \emph{all} models, including those with no
access to the sentiment feature. This is consistent with the asymmetric market response to adverse
news that the behavioural finance literature documents: negative news produces faster, larger and
less predictable moves, so sessions dominated by negative sentiment are harder to forecast for any
model. That the baselines share the pattern confirms it is a property of the market rather than an
artefact of the sentiment pipeline.

Second, and more usefully, the sentiment-augmented model varies \emph{least} across the three
regimes. This is a stability property rather than an accuracy property, and it is the foundation
of the economic result reported in Section~\ref{sec:llm-economic}. A model whose accuracy is
uniform across conditions is more valuable in deployment than one with the same average accuracy
concentrated in favourable conditions, because the sessions in which it degrades are exactly those
in which a position is most exposed.

The same stability appears through time. Tracked across the evaluation period, the
sentiment-augmented model's accuracy is more stable through regime transitions, while the
technical-only baselines fluctuate more sharply. This is the behaviour that adaptive fusion
schemes are designed to produce and that Chapter~\ref{ch:dynfusion} formalises: multi-source
integration provides stability against the degradation of any single information channel.

\section{Confidence Calibration and Economic Significance}
\label{sec:llm-economic}

\subsection{Prediction confidence}

The sentiment-augmented model produces the highest mean prediction confidence of the four models
examined. Taken alone this would be unremarkable and possibly a warning sign --- an overconfident
model is worse than a diffident one. Taken together with comparable accuracy, it means the model
is better \emph{calibrated}: its high-confidence predictions are more often correct and its
low-confidence predictions are correctly flagged as uncertain. This is what feature diversity is
expected to produce, and it is the mechanism by which a statistically unremarkable model can
support better decisions.

\subsection{Simulated trading}

Converting the predictions into positions over a thirty-session simulated trading window, the
sentiment-augmented model produces higher cumulative returns than the baselines despite
comparable statistical accuracy, and avoids large drawdowns during the volatile stretches of the
window.

The divergence between statistical and economic performance is itself the finding of this section,
and it is worth stating explicitly because it recurs throughout the thesis. A model that is right
no more often, but is right \emph{with better-calibrated confidence}, produces better decisions,
because position sizing responds to confidence. A well-calibrated model takes larger positions
when it is more likely to be right and smaller positions when it is not, and the resulting return
stream can differ substantially from that of an equally accurate but poorly calibrated model.

\subsection{Signal-level performance}

Disaggregating by signal direction sharpens the picture. The proposed model performs strongly on
buy signals, with success rates near 72 and 65 per cent against the baselines. Performance on sell
signals is not materially differentiated.

The asymmetry has a plausible mechanism. Positive sentiment amplifies upward price momentum ---
favourable coverage attracts attention, attention attracts participation, participation sustains
the move --- whereas negative sentiment coincides with abrupt declines whose timing is far less
predictable. The framework is therefore differentially useful: it identifies conditions favourable
for entry more reliably than it identifies conditions favourable for exit. A practitioner should
read this as a statement about where the tool applies, not as a general endorsement.

\section{Sentiment Distribution and Lead--Lag Structure}
\label{sec:llm-distribution}

Two diagnostics characterise the sentiment channel itself, independently of any model fitted on
it. They are reported here and revisited in Chapter~\ref{ch:causal}, where sentiment noise is the
explicit object of study.

\textbf{The distribution.} The empirical distribution of the extracted sentiment scores is
approximately symmetric about neutral with moderate tails. This matters as a validity check on the
pipeline. A pipeline that mapped most items to a single class, or that produced a strongly skewed
distribution, would be measuring something other than sentiment --- a source bias, a template
effect, or a failure of the encoder on the domain. A symmetric distribution with genuine
dispersion indicates that the continuous scoring is capturing meaningful variation rather than
producing a biased categorical output, and it justifies the design decision of
Section~\ref{sec:llm-rationale} to retain intensity.

\textbf{Lead--lag structure.} Overlaying the sentiment series on the price series over the
evaluation window shows that strong sentiment spikes frequently \emph{precede} major price
changes. The relationship between magnitude and direction is not simple, however, and two distinct
patterns appear. Sustained periods of positive sentiment accompany gradual price advances, whereas
sharp negative sentiment spikes accompany abrupt declines. The two channels therefore operate on
different timescales: favourable sentiment works through slow accumulation, adverse sentiment
through discrete shocks.

This partially confirms the timeline of sentiment incorporation described in the literature and
partially complicates it. A simple causal ordering --- sentiment moves, price follows --- is not
what the data show. What they show is two regimes of incorporation with different dynamics, which
is one reason a multi-granularity temporal encoder is worth building, as
Chapter~\ref{ch:causal} does.

\section{The Limit of the Result}
\label{sec:llm-limit}

Receiver operating characteristic analysis places \emph{all} models close to the random classifier
line, with areas under the curve clustered near 0.50. The clustering is tight: no model separates
meaningfully from any other, and none separates meaningfully from chance.

This is reported as prominently as the positive findings, and three things follow from it.

\textbf{Language-model sentiment does not overturn market efficiency at this horizon.} At
single-name daily granularity the prediction problem remains close to the efficiency limit of
Fama~\cite{fama1970}, and this thesis does not claim otherwise. The measured increment is of the
order the empirical asset pricing literature leads one to expect~\cite{gu2020}, and a study
reporting a large increment at this horizon would warrant scepticism about its protocol rather
than enthusiasm about its method.

\textbf{The correct baseline is the random classifier, not zero.} A directional accuracy of 47.72
per cent sounds poor until it is placed against a class balance in which upward sessions occur
42.46 per cent of the time --- the precision of the degenerate support vector machine reveals this
directly. Reporting a gain without its baseline is one of the recurring evaluation weaknesses
identified in Section~3.2.2, and this chapter therefore reports the baseline mean of 46.43 per cent
alongside every claim about improvement.

\textbf{The bounded increment is informative, not merely disappointing.} If the value of a
language-model sentiment channel were large at a single-name daily horizon, the correct research
programme would be to refine the channel. Since it is small, the correct programme is to ask
whether value can be created elsewhere in the pipeline. Section~\ref{sec:llm-implication} takes up
that question and Chapter~\ref{ch:allocation} answers it.

\section{Architectural Feasibility at a Real-Time Horizon}
\label{sec:llm-feasibility}

The objective refers to optimising deep learning architectures for real-time prediction, and the
feasibility evidence comes from the optimised architecture of Chapter~\ref{ch:fusion}. That model
has approximately 2.1 million trainable parameters, converges in under two hours on a single
consumer-grade accelerator, and executes a forward pass in milliseconds. Inference cost is
therefore negligible relative to a daily decision horizon.

The binding constraint on ``real time'' in this domain is not inference but ingestion. Retrieving
and scoring hundreds of news items and social posts through transformer encoders introduces
latency measured in seconds to minutes, depending on feed access and batching. For a daily
decision this is comfortable. For an intraday or high-frequency decision it is not, and no
architectural optimisation within the model changes that, because the cost lies in acquiring and
reading the text rather than in predicting from it. The honest statement of scope is therefore
that the framework is real-time at a daily horizon and is not a high-frequency system.

\section{What This Chapter Establishes, and What It Changes}
\label{sec:llm-implication}

Four findings stand, and they must be read together rather than selectively.

\begin{enumerate}
\item Language-model sentiment is the \textbf{highest-gain feature} in the fused space, ahead of
established technical indicators.
\item It is \textbf{orthogonal to price} --- correlation 0.04 with price level, 0.74 with
momentum --- which is the condition under which fusion can add value.
\item It improves \textbf{confidence calibration and simulated economic performance}, with
markedly stronger buy-side signal quality, and it stabilises accuracy across sentiment regimes.
\item It delivers only a \textbf{bounded increment in raw directional accuracy}, 47.72 per cent
against a 46.43 per cent baseline mean, with every model near the random classifier.
\end{enumerate}

The fourth finding is the one that changes the thesis. If the signal channel is weak at the
granularity and horizon where decisions are taken, then refining the signal channel has limited
expected return, and the value of the system must be sought elsewhere. Two candidate explanations
for the weakness present themselves, and they have opposite implications. Either the
\emph{measurement} is inadequate --- the sentiment construct is noisy, and a better extraction
pipeline would recover a larger edge --- or the measurement is sound and the \emph{predictability}
simply is not there at this horizon.

These are distinguishable empirically, and distinguishing them is the task of
Chapter~\ref{ch:causal}, which validates a sentiment construct against four independently
published series and separately measures whether the validated construct forecasts anything.
Chapter~\ref{ch:allocation} then acts on the answer by moving the search for value from the
prediction layer to the decision layer.

\section{Summary}

This chapter isolated a single information source --- news text read by a transformer language
model --- and measured what it contributes to single-name daily directional prediction when fused
with fifteen conventional technical indicators and evaluated against technical-only baselines
under a like protocol.

The sentiment score is the highest-gain feature in the fused space, ranking above established
technical indicators; it is orthogonal to price, correlating 0.04 with price level and 0.74 with
momentum; it produces the highest prediction confidence of the models examined, the most stable
accuracy across sentiment regimes, and the strongest buy-side signal quality, translating into
higher simulated cumulative returns over a thirty-session window. The sentiment-augmented model
attains 47.72 per cent directional accuracy against a baseline mean of 46.43 per cent, a relative
improvement of 2.77 per cent, and does not attain the highest raw accuracy in the comparison.
Receiver operating characteristic analysis places every model close to the random classifier.

The framework runs to completion in a freely accessible hosted environment, so the result is
independently reproducible, and the architectural evidence of Chapter~\ref{ch:fusion} establishes
feasibility at a daily decision horizon with negligible inference cost. The bounded nature of the
accuracy increment is reported rather than concealed, and it is what redirects the remainder of
this thesis from the prediction layer to the decision layer.
